\section*{Chapter \ref{ch:b1:elimination} --- \texorpdfstring{$\beta$}{Beta}-Elimination}

\begin{solution}{exo:b1:elimination:1}
2-Bromobutane: but-1-ene, \cip{E}- and \cip{Z}-but-2-ene; major
\cip{E}-but-2-ene. 2-Bromo-2-methylbutane: 2-methylbut-2-ene (major,
trisubstituted) and 2-methylbut-1-ene. Bromocyclohexane: cyclohexene only.
\end{solution}

\begin{solution}{exo:b1:elimination:2}
E1: $v = k[\mathrm{RBr}]$; E2: $v = k[\mathrm{RBr}][\ce{EtO-}]$. Measure the
\omterm{def:b1:rate-laws:initial-rate}{initial rate} of alkene formation at several ethoxide concentrations: it
grows in proportion for E2, it does not change for E1.
\end{solution}

\begin{solution}{exo:b1:elimination:3}
Along C2--C3, the Br on C2 can be anti to either hydrogen of C3 or to
the C4 methyl group: the two \omterm{def:b1:stereochemistry-in-depth:isomers}{conformations} with an H anti to Br can
react by E2 (giving the \cip{E} and \cip{Z} alkenes); the third cannot
eliminate towards C3.
\end{solution}

\begin{solution}{exo:b1:elimination:4}
SN2 (primary, hydroxide, low temperature); E2 (tertiary, strong bulky
base); SN2 (iodide, a good \omterm{def:b1:electronic-effects:nucleophile}{nucleophile} and \omterm{def:b1:predominant-reaction:strong-acid}{weak base}, \omterm{def:b1:intermolecular-forces-solvents:solvent-classes}{aprotic solvent});
SN1 and E1, E1 increasing with heating.
\end{solution}

\begin{solution}{exo:b1:elimination:5}
C3 carries a single hydrogen. In each diastereomer, placing that H anti to
the Br of C2 fixes the positions of the other groups: the methyl of C2
and the ethyl of C3 are on the same side in one diastereomer and on
opposite sides in the other. One gives \cip{E}-, the other
\cip{Z}-3-methylpent-2-ene: a \omterm{def:b1:nucleophilic-substitution:stereospecific}{stereospecific reaction}.
\end{solution}

\begin{solution}{exo:b1:elimination:6}
The OH is protonated, \ce{R-OH2+}; water leaves, giving the tertiary
\omterm{def:b1:electronic-effects:intermediates}{carbocation} \ce{(CH3)2C+-CH2CH3}; a base (water, \ce{HSO4-}) removes a
proton from a neighbouring carbon. Major product (Zaitsev):
2-methylbut-2-ene.
\end{solution}

\begin{solution}{exo:b1:elimination:7}
The tert-butyl group locks the \omterm{def:b1:stereochemistry-in-depth:conformations}{chair} with itself \omterm{def:b1:stereochemistry-in-depth:conformations}{equatorial}. In the
\textit{trans} isomer the bromine is then \omterm{def:b1:stereochemistry-in-depth:conformations}{equatorial} too, never anti to a
C--H: E2 must wait for the very unfavourable flipped \omterm{def:b1:stereochemistry-in-depth:conformations}{chair}. In the
\textit{cis} isomer the bromine is \omterm{def:b1:stereochemistry-in-depth:conformations}{axial}, with two trans-diaxial
hydrogens: fast E2.
\end{solution}

\begin{solution}{exo:b1:elimination:8}
Ethoxide, small, removes the hydrogen that gives the more substituted
alkene (Zaitsev). tert-Butoxide, bulky, reaches the exposed hydrogens of
the methyl group more easily than the hydrogen of the internal
\ce{CH2}: the less substituted alkene.
\end{solution}

\begin{solution}{exo:b1:elimination:9}
Bromomethane has no $\beta$ carbon. In 1-bromo-2,2-dimethylpropane the
$\beta$ carbon is quaternary and carries no hydrogen.
\end{solution}

\begin{solution}{exo:b1:elimination:10}
Both products need the \omterm{def:b1:electronic-effects:intermediates}{carbocation}, formed in the slow step at the rate
$k_1[\mathrm{RBr}]$. It is then shared between capture by ethanol (\omterm{def:b1:rate-laws:order}{rate
constant} $k_S$) and loss of a proton (\omterm{def:b1:rate-laws:order}{rate constant} $k_E$), both
pseudo-first-order in the solvent: the fraction of alkene is
$k_E/(k_E + k_S)$, independent of the \omterm{def:b1:nucleophilic-substitution:substitution}{substrate} concentration.
\end{solution}

\begin{solution}{exo:b1:elimination:11}
Starting from the \omterm{def:b1:stereochemistry-in-depth:isomers}{conformation} of the \omterm{def:b1:stereochemistry-in-depth:enantiomers}{meso compound} with the two Br anti
(the two phenyls then anti, the two H anti), rotate C1 by $120^\circ$ to
put its H anti to the Br of C2. The two phenyl groups are then on the
same side of the H--Br axis: they end \textit{cis}. On C1, Br outranks
Ph; on C2, Ph outranks H; Br and the phenyl of C2 are \textit{trans}:
\cip{E}-1-bromo-1,2-diphenylethene.
\end{solution}

\begin{solution}{exo:b1:elimination:12}
Only the \omterm{def:b1:stereochemistry-in-depth:conformations}{chair} with the \omterm{def:b1:electronic-effects:nucleophile}{leaving group} \omterm{def:b1:stereochemistry-in-depth:conformations}{axial} reacts: $v = k\,x\,[\mathrm{RY}][\mathrm{B}]$,
with $x$ the fraction of that \omterm{def:b1:stereochemistry-in-depth:conformations}{chair}. Each \omterm{def:b1:stereochemistry-in-depth:conformations}{axial} alkyl group it requires
costs several kJ/mol (\qty{5.7}{kJ/mol} for a methyl), dividing $x$ by about
ten per group at room temperature: with two or three such groups, $x$ falls
to $10^{-2}$ or $10^{-3}$ and the reaction is that much slower.
\end{solution}

\begin{solution}{pb:b1:elimination:1}
\textbf{1.} A \omterm{def:b1:stereochemistry-in-depth:conformations}{chair} with Cl, isopropyl and methyl all \omterm{def:b1:stereochemistry-in-depth:conformations}{equatorial}.
\textbf{2.} Isopropyl and methyl \omterm{def:b1:stereochemistry-in-depth:conformations}{equatorial}, Cl \omterm{def:b1:stereochemistry-in-depth:conformations}{axial}.
\textbf{3.} \omterm{def:b1:stereochemistry-in-depth:enantiomers}{Diastereomers}: they differ only at C1.
\textbf{4.} \omterm{def:b1:electronic-effects:nucleophile}{Leaving group} and $\beta$ hydrogen trans-diaxial.
\textbf{5.} Only in the flipped \omterm{def:b1:stereochemistry-in-depth:conformations}{chair}, with Cl, isopropyl and methyl all
\omterm{def:b1:stereochemistry-in-depth:conformations}{axial}.
\textbf{6.} All three groups \omterm{def:b1:stereochemistry-in-depth:conformations}{axial}: the methyl alone costs about
\qty{5.7}{kJ/mol}, the larger isopropyl group more, the chlorine a little:
together some \qty{13}{kJ/mol}, so a fraction of about
$\exp(-13000/2479) = \num{5e-3}$, the value the problem takes.
\textbf{7.} C2 (one H, it carries the isopropyl group) and C6 (two H).
\textbf{8.} With Cl \omterm{def:b1:stereochemistry-in-depth:conformations}{axial} on C1, the \omterm{def:b1:stereochemistry-in-depth:conformations}{axial} H of C2 (the isopropyl being
\omterm{def:b1:stereochemistry-in-depth:conformations}{equatorial}) and the \omterm{def:b1:stereochemistry-in-depth:conformations}{axial} H of C6: two.
\textbf{9.} In the triaxial \omterm{def:b1:stereochemistry-in-depth:conformations}{chair} the isopropyl group occupies the axial
position of C2, so its H is \omterm{def:b1:stereochemistry-in-depth:conformations}{equatorial}: not anti. C6 has an \omterm{def:b1:stereochemistry-in-depth:conformations}{axial} H: anti.
\textbf{10.} Neomenthyl chloride two, menthyl chloride one.
\textbf{11.} Along C1--C6: Cl (\omterm{def:b1:stereochemistry-in-depth:conformations}{axial}, up) on C1 opposite the \omterm{def:b1:stereochemistry-in-depth:conformations}{axial} H (down)
on C6; the ring bonds and the \omterm{def:b1:stereochemistry-in-depth:conformations}{equatorial} H staggered between.
\textbf{12.} On a \omterm{def:b1:stereochemistry-in-depth:conformations}{chair} no C--H bond is \omterm{def:b1:elimination:periplanar}{syn-periplanar} to the C--Cl bond;
forcing it would mean an eclipsed, strained ring.
\textbf{13.} Removing H of C2: 4-methyl-1-(propan-2-yl)cyclohex-1-ene
(trisubstituted); removing H of C6: 3-methyl-6-(propan-2-yl)cyclohex-1-ene
(disubstituted).
\textbf{14.} The trisubstituted alkene, the more substituted and stable.
\textbf{15.} Both alkenes, \qty{78}{\percent} trisubstituted: consistent.
\textbf{16.} Only the disubstituted alkene: the opposite of Zaitsev's
prediction, imposed by the only anti hydrogen available.
\textbf{17.} Both are stereospecific (anti requirement); neomenthyl chloride
reacts regioselectively (78/22); menthyl chloride gives a single
constitutional isomer.
\textbf{18.} The carbon bearing Cl is secondary and flanked by an isopropyl
group: crowded for SN2, while ethoxide is a \omterm{def:b1:predominant-reaction:strong-acid}{strong base}.
\textbf{19.} $v \propto x\,n_{\mathrm{H}}$: fraction of reactive \omterm{def:b1:stereochemistry-in-depth:conformations}{chair} times
number of anti hydrogens.
\textbf{20.} $0.95 \times 2 = 1.90$.
\textbf{21.} $\num{5.0e-3} \times 1 = \num{5.0e-3}$.
\textbf{22.} \textbf{$k(\text{neomenthyl})/k(\text{menthyl}) = 1.90/\num{5.0e-3}
= 380$}.
\end{solution}
